\section{Условие}
{\bfseries Цель работы:} приобретение практических навыков управления потоками операционной
системы и обеспечения синхронизации между потоками при параллельной обработке данных.

{\bfseries Задание:} составить программу на языке Си, обрабатывающую данные в многопоточном
режиме с использованием штатных средств создания потоков ОС. Ограничение максимального
количества потоков, работающих в один момент времени, задаётся ключом запуска программы.

{\bfseries Вариант 11:} наложить K раз медианный фильтр на матрицу, состоящую из целых чисел.
Размер окна задаётся пользователем.

{\bfseries Требования к решению:}
\begin{itemize}
  \item максимальное число одновременно работающих вычислительных потоков задаётся ключом
    запуска \texttt{--threads P} (лимит общий, включая главный поток);
  \item уметь продемонстрировать количество потоков, используемых программой, штатными
    средствами ОС;
  \item в отчёте привести исследование зависимости ускорения и эффективности алгоритма от
    входных данных и количества потоков с объяснением результатов.
\end{itemize}

{\bfseries Реализация:} Linux/POSIX --- \texttt{pthread\_create}, \texttt{pthread\_join},
\texttt{pthread\_mutex\_*}, \texttt{pthread\_barrier\_t},
\texttt{clock\_gettime(CLOCK\_MONOTONIC)}; сборка --- CMake (на Windows/macOS конфигурация
прерывается с ошибкой).
